\chapter{Discussion and Limitations}
\label{ch:discussion}

\section{What shared anchors solve}

Shared anchors preserve the relative coordinates between a local Gaussian
map and the cameras used to refine it. The controlled pose-correction test
supports this relation, and trajectory accuracy remains close to the
MASt3R-SLAM control. This separation allows appearance refinement to run in
another process without becoming a second pose optimiser.

The construction does not merge neighbouring local maps. Two anchors can
still predict different depths, colours, or extents for the same surface.
Seams and cross-frame colour changes can therefore remain after loop closure.
Joint optimisation of poses, geometry, and appearance is a possible extension.

\section{Map size}
\label{sec:lim:compact}

Dense prediction creates substantially more primitives than the compared
systems. On Replica, our count is about $22.5$--$42.5$ times Photo-SLAM's.
On TUM desk, the ratios are about $66$ to Photo-SLAM and $104$ to MonoGS.
The refiner-off desk run uses $16.4$ times Photo-SLAM's peak GPU memory.

\Cref{tab:curve} keeps the highest-opacity primitives at smaller budgets and
re-renders without optimisation.

\begin{table}[t]
\centering\small
\setlength{\tabcolsep}{4pt}
\begin{tabular}{@{}lrrrr@{}}
\toprule
Gaussians kept & office0 & office1 & room0 & room2 \\
\midrule
full ($\sim$2--3M)     & 26.29 & 22.08 & 25.44 & 23.62 \\
1{,}000{,}000          & 24.68 & 21.71 & 25.01 & 22.87 \\
300{,}000              & 20.23 & 18.96 & 16.26 & 17.87 \\
83{,}372               & 13.54 & 13.97 & \phantom{0}8.68 & 10.43 \\
23{,}019               & 12.04 & 12.08 & \phantom{0}5.16 & \phantom{0}7.74 \\
\bottomrule
\end{tabular}
\caption[PSNR versus map budget, four Replica scenes]{Separate recorded
truncation study on our maps. Highest-opacity primitives are retained
without re-optimisation. Quality remains closer to the full map around one
million primitives and declines sharply at smaller budgets. These points
are not an estimate of optimal quality when trained or refined at each budget.}
\label{tab:curve}
\end{table}


Quality falls sharply below about one million primitives in this test.
The result shows that opacity-ranked removal after full-map optimisation is
insufficient. It does not measure a map trained under the smaller budget.
Budget-aware insertion, merging, pruning, and refinement are therefore
important next steps.

\section{Online refinement}

Moving refinement to a second GPU keeps median tracking latency close to the
refiner-off control on TUM desk. The broader campaign still uses up to
$41.1$~GB on GPU~0 and $14.0$~GB on GPU~1. Seven of nine maps improve LPIPS
before polish, while two larger TUM maps worsen.

The iteration study explains part of this gap. Most of the TUM desk gain
appears by 500 steps, but maps in the causal campaign receive only
$104$--$190$ steps before sequence end. A fixed compute budget supplies fewer
updates per region as the map grows. Adaptive region and view scheduling
should therefore be evaluated together with latency and memory.

\section{Evaluation scope}

The Replica comparison covers eight scenes; the three-method TUM rendering
comparison covers desk. All maps are rendered at common cameras, but input
preprocessing, map size, and total optimisation differ. Photo-SLAM runs to
its own shutdown schedule on Replica. On TUM, our $300$\,s polish is close
to MonoGS's $339$\,s colour phase, but does not equal total compute.

The scoring set excludes keyframes. Other tracked frames may still supervise
our refiner, so this is non-keyframe reconstruction rather than strict
unseen-view evaluation. A stronger generalisation experiment would exclude
scoring frames from every optimisation and model-selection path.

\section{Statistical coverage}

Several reported margins come from one run per configuration. The large
Replica PSNR difference is consistent across all eight scenes, while the
TUM PSNR difference to MonoGS is only $0.035$\,dB and the LPIPS ordering is
reversed. Confidence-ranked and uniform opacity attenuation are also close.
Repeated runs, more sequences, and matched resource budgets are needed to
resolve these smaller effects.
